% Chapter 2, typeset from lectures/L02: the README is the opener and the closing review, and each
% part of Appendix A is a section, in the same order and with the same subsections. Appendix C,
% the guide to running testbenches, is one section after them, with a subsection per part, and
% Appendix B, the exercises, is the last section.

\chapter{Larger Networks, Multiplexers and Submodules}
\label{c2:ch}

{\large\itshape Minimizing a network, the first constructs that need a process, and the first design
built from more than one entity.\par}

\begin{chapterbox}{In this chapter}
\begin{itemize}
  \item \code{process} and \code{case}: sequential statements inside otherwise parallel hardware.
  \item \code{std_logic_vector}: bundling several wires under one name, indexing and slicing.
  \item Multiplexers: a 4-to-1 mux, the smallest thing worth writing a \code{case} for.
  \item Karnaugh maps: why reading sum-of-products off a truth table rarely gives the smallest
    network.
  \item Submodules: writing a block once and instantiating it more than once.
  \item \code{hex_display} on a DE0-CV: eight switches in, two hexadecimal digits out.
  \item Multiplexers at other sizes, in \secref{c2:a4} and \secref{c2:a5}, and the 4-to-1
    re-derived by hand in Exercises~\ref{c2:ex:4} to~\ref{c2:ex:6}.
  \item Running a self-checking testbench with GHDL, the guide every later chapter points back to,
    in \secref{c2:app:c}.
\end{itemize}
\end{chapterbox}

The course this book comes from builds this chapter live, in a particular order. First a 4-to-1
multiplexer in VHDL: four inputs, a two-bit selector, four branches, the smallest thing worth
writing a \code{case} for, and the place where \code{process}, \code{case} and vector indexing all
arrive. Then a Karnaugh map, worked end to end in CircuitVerse: \secref{c2:a1}'s three-variable
example, five \code{1} rows off the truth table, two groups on the map, and \code{X = AB + C} built
as two gates. It is small enough to derive and check by hand, and Exercise~\ref{c2:ex:9} scales the
same technique up to a segment decoder. Then \code{display}, a 7-segment decoder, which is the
multiplexer's \code{case} with sixteen branches instead of four, the only new thing being the
segment table; and \code{hex_display}, two instances of \code{display} with eight switches sliced
into two halves. Finally the design goes onto the board: set the switches, read the hex digits.

Multiplexers at other sizes are for reading rather than presenting, in \secref{c2:a4} and
\secref{c2:a5}. Exercise~\ref{c2:ex:5} builds the 2-to-1, and Exercises~\ref{c2:ex:4}
and~\ref{c2:ex:6} have you re-derive the 4-to-1, by hand and then in VHDL, rather than copy it.

Read the sections that follow before you start the exercises, and install GHDL, following
\secref{c2:app:c}, if you have not already: Chapter~\ref{c1:ch}'s exercises needed it too, and
every exercise from here on that asks you to write a VHDL module does. The board work is shown in
Appendix~\ref{app:quartus} for anyone who has a DE0-CV; no exercise needs one.

% Sections 2.1 to 2.7, from lectures/L02/appendix/a_larger_networks.md: one section per part,
% A.1 to A.7.

\section{Karnaugh Maps: From Truth Table to Minimized Gate Network}
\label{c2:a1}
Reading a sum-of-products equation straight off a truth table, the way \secref{c1:a2} did, always
works but rarely gives the smallest network: every \code{1} row becomes its own AND-term, even when
several rows share most of their inputs and could be combined.

A \textbf{Karnaugh map} (K-map) makes those shared patterns visible by hand, for functions of 2 to 4
variables, with no formal Boolean-algebra manipulation. One worked example follows; further practice
is in \secref{c2:sec:exercises}, and classical minimization theory (don't-cares, Quine--McCluskey)
is out of scope.

The idea is to lay the truth table out as a grid rather than a list, ordering the header labels so
every cell differs from each horizontal and vertical neighbour in exactly one input value, including
around the left/right and top/bottom edges. That ordering is \textbf{Gray code}: for two bits,
\code{00, 01, 11, 10}, where only one bit changes at each step, including from the last entry back
to the first.

Then:
\begin{enumerate}
  \item Write a \code{1} in every cell where the output is \code{1}, and leave the rest blank.
  \item Draw rectangular groups around the \code{1}s:
    \begin{itemize}
      \item Every group must be a rectangle whose width and height are each a power of 2 (1, 2 or
        4), so its size is a power of 2 as well.
        \begin{itemize}
          \item a group of 4 is therefore either a $1\times4$ line or a $2\times2$ block, and both
            are allowed.
          \item cells that touch only at a corner are not a rectangle, so they cannot be grouped.
        \end{itemize}
      \item Groups may overlap, and may wrap around an edge of the grid.
    \end{itemize}
  \item Make every group as large as possible, and cover every \code{1} at least once. A group of 1
    keeps all its inputs in the term, a group of 2 drops one, a group of 4 drops two, and so on.
    \begin{itemize}
      \item Then use as few groups as will do it: drop any group whose cells are every one of them
        already covered by another. Nothing above forbids a redundant group, and a redundant group
        is a redundant AND-term in the answer, correct but not minimal.
    \end{itemize}
  \item For each group, keep only the input(s) that hold the same value in every cell of that
    group, dropping the rest, and write that as one AND-term:
    \begin{itemize}
      \item write a kept input directly if that constant value is \code{1}, and primed if it is
        \code{0}.
      \item so a group covering only cells where \code{B = 0} contributes the term \code{B'}.
    \end{itemize}
  \item \code{X} is the OR of all the group terms.
\end{enumerate}

\subsection{Worked example}
\label{c2:kmapexample}
Three inputs \code{A}, \code{B}, \code{C}, one output \code{X}:

\begin{narrowtable}{@{}cccc@{}}
  \thead{A} & \thead{B} & \thead{C} & \thead{X} \\
  \midrule
  0 & 0 & 0 & 0 \\
  0 & 0 & 1 & 1 \\
  0 & 1 & 0 & 0 \\
  0 & 1 & 1 & 1 \\
  1 & 0 & 0 & 0 \\
  1 & 0 & 1 & 1 \\
  1 & 1 & 0 & 1 \\
  1 & 1 & 1 & 1 \\
\end{narrowtable}

\noindent Lay it out with \code{AB} in Gray code down the rows and \code{C} across the columns,
filling in a \code{1} wherever \code{X = 1} (Figure~\ref{c2:fig:karnaugh_filled}).

\bookfigure{karnaugh_filled}{The Karnaugh map for \code{X}, filled in from the truth table:
\code{AB} in Gray code down the rows, \code{C} across the columns.}{c2:fig:karnaugh_filled}

\noindent The entire \code{C = 1} column is four \code{1}s, a group of 4. Every cell shares
\code{C = 1} and nothing else is constant (\code{A} and \code{B} both vary), so it reduces to the
single-variable term \code{C}:

\bookfigure{karnaugh_group_c}{The same map with the four-cell \code{C = 1} group marked, which
reduces to the term \code{C}.}{c2:fig:karnaugh_group_c}

\noindent One \code{1} is still uncovered, on the \code{AB = 11} row. Group it with its neighbour to
the right, which the \code{C} group already covers. Every cell in the pair shares \code{A = 1} and
\code{B = 1}, and \code{C} varies, so it reduces to \code{AB}:

\bookfigure{karnaugh_groups}{The same map with the two-cell \code{AB = 11} group added, overlapping
the \code{C} group in one cell.}{c2:fig:karnaugh_groups}

\noindent Every \code{1} is now covered; the \code{AB = 11, C = 1} cell is in both groups, and
overlap is fine. Summing gives:

\begin{console}
X = AB + C
\end{console}

\noindent Reading the truth table directly, as Chapter~\ref{c1:ch} did, would have given five
AND-terms, one per \code{1} row. The K-map found an equivalent two-term equation, which is the entire
value of doing this by hand.

\subsection{Realizing the network}
\label{c2:kmapnetwork}
\code{X = AB + C} is two gates: an \code{AND} computing \code{AB}, feeding an \code{OR} whose other
input is \code{C}.

\bookimage{0.55}{lectures/L02/appendix/images/net.png}{The gate network for the minimized
expression \code{X = AB + C}: an \code{AND} gate feeding an \code{OR}, built in
CircuitVerse.}{c2:fig:net}

\noindent The course builds this network by hand in CircuitVerse, live; build it yourself exactly as
you built \code{or_gate}'s in Chapter~\ref{c1:ch}. \Secref{c2:a2} translates it into VHDL.

\section{Translating the Karnaugh-Map Network into VHDL}
\label{c2:a2}
\code{X = AB + C} translates straight into VHDL. Three inputs, one output:

\begin{vhdlcode}
library ieee;
use ieee.std_logic_1164.all;

entity gate_network is
    port(a, b, c: in std_logic;
         x      : out std_logic);
end entity;
\end{vhdlcode}

\noindent The architecture uses an internal signal \code{ab} to hold the \code{AND} gate's output,
mirroring the two-gate network one-for-one:

\needspace{12\baselineskip}
\begin{vhdlcode}
architecture behaviour of gate_network is
signal ab: std_logic;
begin
    -- AND gate: AB.
    ab <= a and b;

    -- OR gate: AB + C.
    x <= ab or c;
end architecture;
\end{vhdlcode}

\noindent Without the intermediate signal this collapses to \code{x <= (a and b) or c;}. Both are
equally valid and synthesize to the same hardware; use whichever is more readable for the network at
hand.

The buildable version, \file{lectures/L02/gate_network/gate_network.vhd}, uses the collapsed form,
since at two gates there is nothing to gain from naming the intermediate. It is printed in full
below; set it beside the version above and satisfy yourself they describe the same circuit. That is
the claim this section is making, and it costs nothing to check.

\vhdlfile{lectures/L02/gate_network/gate_network.vhd}

To run it on hardware, assign the ports to switches and an LED with the DE0-CV's pin planner
(Appendix~\ref{app:quartus}). Nothing here depends on specific pin numbers.

\section{\code{std_logic_vector}: More Than One Wire}
\label{c2:a3}
Everything so far has been \code{std_logic}, one wire carrying one bit. Bundling several wires under
one name is what \code{std_logic_vector} is for.

\begin{vhdlcode}
signal count: std_logic_vector(3 downto 0);
\end{vhdlcode}

\noindent That declares four wires: \code{count(3)} down to \code{count(0)}. The range is
\code{3 downto 0} rather than \code{0 to 3} so the leftmost bit is the most significant, the way you
would write the number on paper. This book uses \code{downto} everywhere, and so does almost all the
VHDL you will meet.

A vector is a port type like any other:

\begin{vhdlcode}
entity display is
    port(number: in  std_logic_vector(3 downto 0);
         hex   : out std_logic_vector(6 downto 0));
end entity;
\end{vhdlcode}

\subsection{Values}
\label{c2:values}
A single bit takes single quotes, a vector double quotes:

\begin{vhdlcode}
x     <= '1';                -- std_logic: one bit
count <= "1010";             -- std_logic_vector(3 downto 0): four bits
count <= (others => '0');    -- every bit '0', whatever the width
\end{vhdlcode}

\noindent The last is an \emph{aggregate}: \code{(others => '0')} means ``all the remaining bits are
\code{'0'}'', and since none were named individually, that is all of them. It keeps working when the
width changes, and it appears throughout the rest of the book.

\subsection{Indexing and slicing}
\label{c2:slicing}
One bit out of a vector is a \code{std_logic}; a contiguous run is a \emph{slice}, itself a vector:

\begin{vhdlcode}
carry     <= count(3);
high_half <= input(7 downto 4);
\end{vhdlcode}

\noindent Slicing is how one wide port feeds several narrower ones, which is what \secref{c2:a6}'s
two-digit display does with its eight switch inputs.

One detail is not what a software engineer expects. \code{input(7 downto 4)} has the index range
\code{7 downto 4}, not \code{3 downto 0}. Connected to a port declared
\code{std_logic_vector(3 downto 0)}, the two are matched \textbf{by position, left to right}, not by
index number: \code{input(7)} drives bit 3, \code{input(6)} drives bit 2, and so on. The lengths
have to match. The numbering does not.

\subsection{A vector is a bundle of bits, not a number}
\label{c2:bundle}
\code{count <= "1010"} does not make \code{count} equal ten. \code{std_logic_vector} says nothing
about what the bits \emph{mean}, so it offers only the operations that make sense for a bundle of
wires: the logical operators applied bit by bit (\code{a and b}, \code{not a}, \code{a xor b}),
comparison (\code{count = "1010"}), and concatenation (\code{"10" & "11"} is \code{"1011"}).

Arithmetic is not on that list. \code{count + 1} will not compile, because nothing has said whether
those four wires are an unsigned number, a signed one, or four unrelated signals.

This book does not work around that with a numeric vector type. Anything counted is a plain integer
with a range:

\begin{vhdlcode}
signal count: natural range 0 to 9;
\end{vhdlcode}

\noindent which you can add to and compare directly, and which synthesis turns into exactly the
flip-flops the range needs. Vectors stay what they are: bundles of wires, for things that really are
bundles of wires, like eight switches or seven display segments.

The two worlds meet only at a boundary, and the conversions appear where they are needed:

\begin{vhdlcode}
count_v <= std_logic_vector(to_unsigned(count, 4));   -- integer to vector
count   <= to_integer(unsigned(count_v));             -- vector to integer
\end{vhdlcode}

\noindent Both come from \code{ieee.numeric_std}, which is why some files add a third \code{use}
clause.

\section{Multiplexers}
\label{c2:a4}
A \textbf{multiplexer} (``mux'') routes exactly one of several data inputs to a single output, chosen
by a separate set of \textbf{selector} inputs. An \code{N}-to-1 mux has \code{N} data inputs and
\code{ceil(log2(N))} selector bits.

The simplest case, a 2-to-1 mux with inputs \code{A}, \code{B}, selector \code{S} and output
\code{X}:

\begin{narrowtable}{@{}cc@{}}
  \thead{S} & \thead{X} \\
  \midrule
  0 & B \\
  1 & A \\
\end{narrowtable}

\noindent As a Boolean equation, \code{X = S'B + SA}: for every input combination exactly one
AND-term is active, and that term's data input passes through.

Larger muxes are the same idea with more selector bits, though \emph{which} data input a given
selector value names is a convention that changes from part to part, and has to be read off the
specification rather than assumed. The 8-to-1 below is the opposite way round from the 2-to-1
above: there \code{0} names the second input, here \code{000} names the first. An 8-to-1 mux with
data inputs \code{A}--\code{H} and selector \code{S[2:0]}:

\begin{narrowtable}{@{}cc@{}}
  \thead{S[2:0]} & \thead{X} \\
  \midrule
  000 & A \\
  001 & B \\
  010 & C \\
  011 & D \\
  100 & E \\
  101 & F \\
  110 & G \\
  111 & H \\
\end{narrowtable}

\noindent Read into sum-of-products one term per row, the \code{000} row contributes
\code{A*S2'S1'S0'} and the rest follow the same shape, each gating its data input with the selector
combination that selects it:

\begin{console}
X = AS2'S1'S0' + BS2'S1'S0 + ... + HS2S1S0
\end{console}

\noindent Worth writing out once by hand, not worth memorizing. The pattern is the point, and it is
the same at any size.

As gates: one 4-input \code{AND} per data input (the data bit plus all three selector bits),
inverting selector bits as needed, with all eight \code{AND} outputs summed by a tree of \code{OR}
gates (Figure~\ref{c2:fig:mux8}).

\bookimage[htbp]{0.85}{lectures/L02/appendix/images/mux_8to1.png}{The 8-to-1 multiplexer as gates: eight
4-input \code{AND} gates, each gating one data input with the selector combination that selects it,
feeding a tree of \code{OR} gates.}{c2:fig:mux8}

\noindent This is a real building block. The ATmega328P (the chip on an Arduino Uno/Nano) uses one
internally so its analog-capable pins, plus an internal bandgap reference and a temperature sensor,
can share a single analog-to-digital converter, selected by the four \code{MUX} bits of its
\code{ADMUX} register.

\section{Multiplexers in VHDL: \code{process} and \code{case}}
\label{c2:a5}
Every statement so far has been \textbf{concurrent}: a standing description of a wire, all active at
once. A multiplexer is the first circuit that is awkward to write that way, and needs two new
constructs:
\begin{itemize}
  \item a \textbf{\code{process}}, a block whose contents run sequentially, top to bottom, like
    ordinary code. It re-runs whenever any signal in its \textbf{sensitivity list} changes, and
    from the architecture's point of view it is still just one more concurrent statement. Every
    process in this book has this shape; Chapter~\ref{c3:ch} reuses it for clocked logic, with the
    clock in the sensitivity list instead of the inputs being decoded.
  \item a \textbf{\code{case} statement}, VHDL's \code{switch}. Like \code{if} and \code{for} it is
    a \emph{sequential} statement, legal only inside a process (or a subprogram, which this book
    does not use), and every value the selector can take must be covered, which
    \code{when others} guarantees.
\end{itemize}
``Pick one of several inputs based on a selector'' \emph{is} a \code{case} statement, almost word for
word. The 2-to-1 mux from \secref{c2:a4}:

\needspace{11\baselineskip}
\begin{vhdlcode}
MUX_PROCESS: process(a, b, sel) is
begin
    case (sel) is
        when '0'    => x <= b;
        when '1'    => x <= a;
        when others => x <= '0';
    end case;
end process;
\end{vhdlcode}

\noindent Two points from that:
\begin{itemize}
  \item Single-bit values take single quotes, \code{'0'}/\code{'1'}, as opposed to the double quotes
    used for vectors like \code{"000"}.
  \item \code{when others} is required rather than defensive. VHDL insists a \code{case} cover every
    value of its selector's type, and \code{std_logic} has nine (\secref{c1:a4}), so \code{'0'} and
    \code{'1'} alone never make one exhaustive. This is a language rule, not a claim about the wire:
    a synthesized net carries a \code{0} or a \code{1} and nothing else, and the other seven values
    are the simulator's, where a signal sitting at \code{'U'} or \code{'X'} is a bug rather than a
    level.
\end{itemize}

\subsection{Worked example: the 8-to-1 mux}
\label{c2:mux8}
Same shape with a 3-bit selector and eight branches, one per row of \secref{c2:a4}'s truth table.
The data inputs \code{A}--\code{H} become the vector \code{inputs(7 downto 0)}, with \code{A} as
\code{inputs(7)} down to \code{H} as \code{inputs(0)}, so selector \code{"000"} routes
\code{inputs(7)} and \code{"111"} routes \code{inputs(0)}.

The buildable module is \file{lectures/L02/mux_8to1/mux_8to1.vhd}:

\vhdlfile{lectures/L02/mux_8to1/mux_8to1.vhd}

Compare it against \secref{c2:a4}'s eight-term equation for the same circuit. The equation describes
the gates; the \code{case} describes the \emph{intent} and lets the synthesis tool produce the gates.
Both are correct and both synthesize to the same hardware. That gap, between describing structure
and describing behaviour, is most of what a hardware description language buys you.

\section{Building a Design from Submodules}
\label{c2:a6}
Everything so far has been one entity with one architecture. That stops being enough the moment a
design needs the same block twice.

The DE0-CV has six 7-segment displays, and you want to drive two, each showing one hexadecimal
digit. The logic turning a 4-bit number into seven segment signals is the same for both. You could
write it twice inside one architecture and nothing would stop you, but the second copy is a
liability: every fix has to be made in both, and the day they drift apart is the day you lose an
evening working out which one is wrong.

Write the block once, as its own entity, and \emph{instantiate} it twice.

\subsection{The submodule}
\label{c2:submodule}
An ordinary module. Nothing about it says ``I am a submodule'':

\begin{vhdlcode}
entity display is
    port(number: in  std_logic_vector(3 downto 0);
         hex   : out std_logic_vector(6 downto 0));
end entity;
\end{vhdlcode}

\noindent Whether it ends up as a whole design or one piece of a larger one is decided by whoever
instantiates it, not by the module.

\subsection{Instantiating it}
\label{c2:instantiate}
With \code{entity work.<name>}, inside another architecture:

\begin{vhdlcode}
entity hex_display is
    port(input     : in  std_logic_vector(7 downto 0);
         hex1, hex0: out std_logic_vector(6 downto 0));
end entity;

architecture structure of hex_display is
begin
    display1: entity work.display
        port map(input(7 downto 4), hex1);

    display0: entity work.display
        port map(input(3 downto 0), hex0);
end architecture;
\end{vhdlcode}

\begin{itemize}
  \item \code{display1} and \code{display0} are \textbf{instance labels}, naming the two copies.
    They must be unique within the architecture, and they are not the module's name: both instances
    are \code{display}.
  \item \code{entity work.display} selects the module. \code{work} is the library your analyzed
    files land in by default, in GHDL and Quartus alike.
  \item \code{port map(...)} connects the instance's ports to signals in the enclosing architecture,
    \textbf{in the order that entity declares them}. \code{display} declares \code{number} then
    \code{hex}, so the first item drives \code{number} and the second is driven by \code{hex}.
  \item \code{input(7 downto 4)} is a \textbf{slice}, which the instance sees as a
    \code{std_logic_vector(3 downto 0)}.
\end{itemize}
Note what \code{hex_display}'s architecture does \emph{not} contain: any logic. It is wiring and
nothing else. Every gate in the finished design comes from the two \code{display} instances.

\subsection{Positional and named association}
\label{c2:association}
The \code{port map} above is \textbf{positional}, matching ports by order. VHDL also allows
\textbf{named} association, which you will meet in other people's code:

\begin{vhdlcode}
display1: entity work.display
    port map(number => input(7 downto 4), hex => hex1);
\end{vhdlcode}

\noindent The two are equivalent. This book uses positional everywhere, which means an entity's port
declaration order is something you have to get right. That is not a formality: it is exactly what
the testbenches you are given rely on, which is why every exercise's self-check line spells out its
ports \emph{in order}.

\subsection{What this costs in hardware}
\label{c2:cost}
Nothing surprising. Instantiating a module twice puts two copies of its gates on the FPGA, exactly
as writing the logic twice would. The saving is in the source, not the silicon: one description, one
place to fix, and a design you can read.

\section{Checking Your Work}
\label{c2:a7}
Every VHDL module you write for an exercise comes with a \textbf{self-checking testbench}: a small
VHDL file that drives your inputs, checks your outputs, and tells you which case failed.
Chapter~\ref{c1:ch}'s Exercise~\ref{c1:ex:2} already shipped one, and \secref{c1:a5} gave you the
three commands that run it.

You do not write testbenches in this book. Running them is enough to verify every exercise on your
own laptop, with no FPGA board.

\Secref{c2:app:c} is the full guide: what a testbench is and why its entity has no ports, the
analyze $\to$ elaborate $\to$ run sequence, how to read a failed assertion, and how your module's
\textbf{port order} has to match what the testbench expects.

Run one now, against this chapter's worked example:

\begin{shell}
cd lectures/L02/gate_network
ghdl -a --std=93 gate_network.vhd gate_network_tb.vhd
ghdl -e --std=93 gate_network_tb
ghdl -r --std=93 gate_network_tb --assert-level=error --stop-time=10ms
\end{shell}

% Section 2.8, from lectures/L02/appendix/c_testbenches.md: the course's guide to running
% testbenches, one subsection per part, C.1 to C.7. Every later chapter points back here.

\section{Verifying Your VHDL with a Testbench}
\label{c2:app:c}
This section is the book's guide to \textbf{running} testbenches, and every later chapter's
exercises point back here. You do not write testbenches in this book, but you run one after every
VHDL exercise, because with no FPGA board of your own that is what verifies your VHDL. Every worked example and every exercise ships a
ready-made \file{<module>_tb.vhd}.

\subsection{What you need}
\label{c2:c1}
\textbf{GHDL}, the open-source VHDL analyzer and simulator. Check with \code{ghdl --version}; install
on WSL/Ubuntu with \code{sudo apt -y install ghdl}.

Every command below uses \code{--std=93}, the VHDL standard this book targets.

\subsection{What a testbench is}
\label{c2:c2}
A testbench is itself a VHDL module, but a special one:
\begin{itemize}
  \item Its entity has \textbf{no ports}: nothing connects to the outside, because it is the top of
    the simulation.
  \item Its architecture \textbf{instantiates your design}, the ``device under test'' labelled
    \code{dut}, and drives its inputs.
  \item It \textbf{checks} each output with an \code{assert}. If an output is wrong, the assertion
    fires with a message and the simulation stops.
\end{itemize}
So instead of you reading a waveform to decide whether the output is right, the testbench decides,
and only speaks up when something is wrong.

\subsection{Running a testbench: analyze \texorpdfstring{$\to$}{->} elaborate
  \texorpdfstring{$\to$}{->} run}
\label{c2:c3}
Three steps. Pass \code{--std=93} to all three, and give the \textbf{testbench entity name}, not a
filename, to the last two:

\begin{shell}
ghdl -a --std=93 <module>.vhd <module>_tb.vhd  # analyze both files (module first)
ghdl -e --std=93 <module>_tb                   # elaborate the testbench
ghdl -r --std=93 <module>_tb --assert-level=error --stop-time=10ms   # run it
\end{shell}

\begin{itemize}
  \item \code{--assert-level=error} makes a failed check stop the run and return a
    \textbf{non-zero exit code}, a real pass/fail you can check with \code{echo $?}. The testbenches
    here raise failures at \code{severity failure}, which stops the run on its own, so you get a
    real pass/fail even without the flag. Keep it on anyway: it is the habit that also works with
    testbenches reporting at a lower severity.
  \item \code{--stop-time=10ms} caps how long the simulation may run. A testbench often waits for
    your module to do something, with a line like \code{wait until tx = '0';}. If your module never
    does, that wait never finishes and without a stop time the simulation hangs: no output, no
    error, just a terminal sitting there. 10\,ms of \emph{simulated} time is far more than any
    testbench here needs, and takes well under a second of real time.
  \item A passing run prints its ``all checks passed'' note and exits 0. Every testbench in this
    book ends with that note, so if you do not see it, the run did not finish: it either failed a
    check or hit the stop time.
\end{itemize}

\subsection{Checking an exercise}
\label{c2:c4}
Each VHDL exercise ships its testbench in its own directory. For example, this chapter's
Exercise~\ref{c2:ex:3} has \file{lectures/L02/exercises/xyz_logic/xyz_logic_tb.vhd}. Write your
module in \textbf{that same directory}, named exactly as the exercise asks, then run the three
commands from there:

\begin{shell}
cd lectures/L02/exercises/xyz_logic
ghdl -a --std=93 xyz_logic.vhd xyz_logic_tb.vhd
ghdl -e --std=93 xyz_logic_tb
ghdl -r --std=93 xyz_logic_tb --assert-level=error --stop-time=10ms
\end{shell}

\noindent Each testbench header lists these exact commands.

The testbench connects to your module \textbf{by position}, not by name, so your entity must declare
its ports in the \textbf{same order} the exercise lists them. The names are yours to choose, but use
the ones specified anyway, so your module and the testbench's error messages talk about the same
signals.

A port declared in the wrong order fails in one of two ways, neither of which says ``wrong order'':
\begin{itemize}
  \item if the misplaced ports have \textbf{different} types or widths, analysis stops with
    \code{can't associate "sel" with port "sel"}, pointing at the testbench's \code{port map} line.
  \item if they have the \textbf{same} type, several \code{std_logic} ports say, it analyzes and
    elaborates quite happily, then fails at run time as an ordinary assertion failure, which reads
    exactly like a logic bug in your design.
\end{itemize}
So when a testbench fails on the very first case it checks, re-read your port order before hunting
through your architecture. The same applies to \textbf{generics}: a module that carries them must
declare exactly the ones the exercise specifies, in order.

\subsection{When the design needs more than one file}
\label{c2:c5}
Some exercises build a module out of another one, so the \code{ghdl -a} line names several files
instead of two. The elaborate and run steps are unchanged: they still take the testbench
\textbf{entity name} only.

\needspace{5\baselineskip}
\begin{shell}
ghdl -a --std=93 <subcomponent>.vhd <module>.vhd <module>_tb.vhd
ghdl -e --std=93 <module>_tb
ghdl -r --std=93 <module>_tb --assert-level=error --stop-time=10ms
\end{shell}

\begin{itemize}
  \item \textbf{Order matters on the analyze line}: name a module before anything that instantiates
    it, so GHDL has already seen the entity when it reaches the \code{entity work.<name>} that
    refers to it. Subcomponents first, your top level next, the testbench last.
  \item No exercise directory ships a subcomponent. Every module a design instantiates is one you
    wrote in an earlier exercise and copy in yourself; the exercise says which, and where from.
  \item The exercises that need this are this chapter's \code{hex_display}
    (Exercise~\ref{c2:ex:8}, which instantiates your \code{display}), Chapter~\ref{c4:ch}'s
    \code{led_toggle_sync2}, Chapter~\ref{c7:ch}'s \code{blinker} (\code{reset_sync} and
    \code{timer}) and \code{walking_led} (those two plus \code{button_sync}), and all of
    Chapter~\ref{c8:ch}'s, where the receiver names five files and the transmitter four.
  \item Every testbench header lists the exact command for its own exercise.
\end{itemize}

\subsection{Reading the result}
\label{c2:c6}
A failed check looks like this:

\begin{console}
xyz_logic_tb.vhd:37:13:@10ns:(assertion failure): xyz_logic: wrong X with a='0' b='0' c='0' d='0', expected '0' but got '1'!
ghdl:error: assertion failed
\end{console}

\noindent Read it in order: \code{file:line} is where the assert is, \code{@time} is when it fired,
then the message. It is the same shape as a GHDL compile error, so it maps onto what you already
know. No failure means every case matched.

\subsection{A complete example}
\label{c2:c7}
A combinational testbench in the shape every testbench in this book uses, here for
Chapter~\ref{c1:ch}'s \code{or_gate}:

\begin{vhdlcode}
library ieee;
use ieee.std_logic_1164.all;
use ieee.numeric_std.all;

entity or_gate_tb is
end entity;

architecture behaviour of or_gate_tb is
signal a, b, x: std_logic;
begin
    dut: entity work.or_gate
        port map(a, b, x);   -- positional: same order as or_gate's port clause.

    SIM_PROCESS: process is
    begin
        for i in 0 to 3 loop
            (a, b) <= std_logic_vector(to_unsigned(i, 2)); -- apply inputs
            wait for 10 ns;                                -- let them settle
            assert x = (a or b)                            -- check AFTER the wait
                report "or_gate: wrong output with a=" & std_logic'image(a)
                     & " b=" & std_logic'image(b)
                     & ", expected " & std_logic'image(a or b)
                     & " but got " & std_logic'image(x) & "!"
                severity failure;
        end loop;
        report "or_gate: all checks passed!" severity note;
        wait;                                              -- end the simulation
    end process;
end architecture;
\end{vhdlcode}

\noindent Two rules keep a combinational testbench correct:
\begin{itemize}
  \item \textbf{Apply $\to$ wait $\to$ check, in that order.} A signal assignment only
    \emph{schedules} a value, so the inputs and output do not update until the \code{wait}.
    Checking before it tests the previous iteration's values.
  \item \textbf{End with \code{wait;}} so the stimulus process suspends and the simulation stops on
    its own.
\end{itemize}
Clocked designs need a little more, a clock-generator process and a way to stop it, which is exactly
the machinery the provided testbenches already contain. For the exercises you only ever write the
module.

% Section 2.9, from lectures/L02/README.md: the objectives, the questions to test yourself, and the
% next lecture.

\needspace{10\baselineskip}
\section{Review}
\label{c2:sec:review}

\subsection*{What you should be able to do now}
After this chapter, you should be able to:
\begin{itemize}
  \item Derive a minimized equation from a Karnaugh map and turn it into a VHDL architecture.
  \item Read and write \code{std_logic_vector} ports and signals.
  \item Write a \code{process} with a correct sensitivity list and a \code{case} covering every input
    value.
  \item Build a design out of more than one entity with a positional \code{port map}, and say why
    instantiating a module twice puts two copies of its gates on the FPGA.
  \item Verify a module by running its testbench, and read a failed assertion.
\end{itemize}

\needspace{6\baselineskip}
\subsection*{Questions to test yourself}
You should be able to explain:
\begin{enumerate}
  \item Given a truth table, can you derive its equation both directly and via a Karnaugh map? What
    is different about the two results?
  \item Why can a \code{case} statement not appear directly in an architecture body?
  \item What does \code{when others} protect you against, given that a selector ``can only'' be
    \code{0} or \code{1}?
  \item In \code{display1: entity work.display}, what does \code{display1} name, and what does
    \code{work.display} name? Could a second instance reuse the same label?
  \item \code{hex_display}'s architecture contains no logic at all. Where do the gates in the finished
    design come from?
  \item You instantiated \code{display} twice rather than writing its \code{case} twice. Does the FPGA
    end up with one copy of that logic or two? So what did you actually gain?
  \item This book always writes \code{port map} positionally. What breaks if an entity's ports are
    declared in a different order from the one its testbench expects, and when would you find out?
\end{enumerate}

\needspace{6\baselineskip}
\subsection*{Looking ahead}
Chapter~\ref{c3:ch} is about sequential logic, circuits that remember:
\begin{itemize}
  \item The D latch, and the D flip-flop built from a pair of them.
  \item Clocking, and when a signal assignment actually takes effect.
  \item Edge detection: a button press that toggles an LED.
\end{itemize}

% Section 2.10, from lectures/L02/appendix/b_exercises.md.

\section{Exercises}
\label{c2:sec:exercises}

\begin{aside}
\textbf{How to work these.} Each exercise is tagged with the kinds of work it asks for:
\kinds{Paper} by hand, \kinds{CircuitVerse} in the browser, \kinds{VHDL} in a module of your own.

\textbf{How to check your work.} Every exercise that asks you to write a VHDL module ships a
self-checking testbench under \file{lectures/L02/exercises} in the companion repository. Write your
module in its \file{exercises/<module>} directory, using the entity name and the \textbf{port order}
the exercise specifies, then run it with GHDL; \secref{c2:app:c} has the three commands.

No FPGA board is needed for any exercise. Your job is to get the VHDL right, and the testbench is
how you confirm it. Short answers to the parts you work on paper are in
Appendix~\ref{app:answers}; reference solutions to the VHDL are published in the companion
repository, under \file{lectures/L02/solutions}. Try each exercise yourself before looking at
either.
\end{aside}

% ------------------------------------------------------------------------------------------
\exerciseset{Karnaugh Maps}
\label{c2:set:1}

\exercise{A three-variable map}{Paper, CircuitVerse}
\label{c2:ex:1}
Derive a minimized equation for \code{X} from the Karnaugh map below, then realize the
corresponding gate network and simulate it in CircuitVerse:

\begin{narrowtable}{@{}cc@{}}
  \thead{ABC} & \thead{X} \\
  \midrule
  000 & 1 \\
  001 & 1 \\
  010 & 0 \\
  011 & 0 \\
  100 & 1 \\
  101 & 1 \\
  110 & 0 \\
  111 & 0 \\
\end{narrowtable}

\note[Tip] Lay \code{AB} down the rows in Gray-code order (\code{00, 01, 11, 10}) and \code{C}
across the columns, exactly as in \secref{c2:a1}'s worked example, before you start grouping
\code{1}s.

\exercise{A four-variable map}{Paper, CircuitVerse}
\label{c2:ex:2}
Derive a minimized equation for \code{X} from the Karnaugh map below, then realize the
corresponding gate network and simulate it in CircuitVerse:

\begin{narrowtable}{@{}cc@{}}
  \thead{ABCD} & \thead{X} \\
  \midrule
  0000 & 0 \\
  0001 & 1 \\
  0010 & 0 \\
  0011 & 1 \\
  0100 & 0 \\
  0101 & 0 \\
  0110 & 0 \\
  0111 & 0 \\
  1000 & 0 \\
  1001 & 1 \\
  1010 & 0 \\
  1011 & 1 \\
  1100 & 1 \\
  1101 & 0 \\
  1110 & 1 \\
  1111 & 0 \\
\end{narrowtable}

\note[Tip] With four variables, put two of them (in Gray code) on each axis, so you get a
$4\times4$ grid. Groups can still wrap around every edge of the grid, not just left-right; remember
to check the top/bottom wraparound too.

% ------------------------------------------------------------------------------------------
\needspace{16\baselineskip}
\exerciseset{Circuits in VHDL}
\label{c2:set:2}

\exercise{A three-output network}{Paper, CircuitVerse, VHDL}
\label{c2:ex:3}
A gate network has:
\begin{itemize}
  \item Four inputs, \code{ABCD}.
  \item Three outputs, \code{XYZ}.
\end{itemize}
It is given by the truth table below:

\begin{narrowtable}{@{}cc@{}}
  \thead{ABCD} & \thead{XYZ} \\
  \midrule
  0000 & 001 \\
  0001 & 010 \\
  0010 & 001 \\
  0011 & 011 \\
  0100 & 100 \\
  0101 & 110 \\
  0110 & 100 \\
  0111 & 111 \\
  1000 & 101 \\
  1001 & 110 \\
  1010 & 101 \\
  1011 & 111 \\
  1100 & 000 \\
  1101 & 010 \\
  1110 & 000 \\
  1111 & 011 \\
\end{narrowtable}

\begin{parts}
  \item Treat \code{X}, \code{Y}, and \code{Z} as three separate single-output functions of
    \code{ABCD}, and derive a minimized equation for each via a Karnaugh map.
  \item Realize the resulting gate network and simulate it in CircuitVerse; confirm all three
    outputs against the truth table above.
  \item Implement the design in VHDL, in a module of your own with four \code{std_logic} inputs and
    three \code{std_logic} outputs (\secref{c1:a4} and \secref{c2:a2} show the pattern for turning a
    derived equation into concurrent VHDL assignments).
  \item Run the testbench (see the aside at the start of these exercises) and confirm all 16 rows
    pass.
\end{parts}
If a row fails, the assertion message names the input combination:
\begin{itemize}
  \item Go back to the Karnaugh map for whichever of \code{X}, \code{Y}, \code{Z} is wrong.
  \item A single wrong output on a single row is nearly always one mis-grouped \code{1} in one map.
\end{itemize}
\note[Tip] Two of the three outputs here turn out to depend on only one or two of the four inputs
once minimized. Don't assume every output needs all four.

\modulefigure{xyz_logic}

\selfcheck name your VHDL entity \code{xyz_logic}, with ports \code{a}, \code{b}, \code{c},
\code{d} (in) and \code{x}, \code{y}, \code{z} (out), declared in that order; its testbench is in
\file{exercises/xyz_logic} and checks all 16 rows of the truth table above.

% ------------------------------------------------------------------------------------------
\exerciseset{Multiplexers}
\label{c2:set:3}

\exercise{The 4-to-1 mux by hand}{Paper, CircuitVerse}
\label{c2:ex:4}
A 4-to-1 multiplexer has:
\begin{itemize}
  \item Four data inputs, \code{A}--\code{D}.
  \item Two selector bits, \code{S[1:0]}.
  \item One output, \code{X}.
\end{itemize}
Take the selector to name its data input directly: \code{S[1:0] = "00"} selects \code{A},
\code{"01"} selects \code{B}, \code{"10"} selects \code{C}, and \code{"11"} selects \code{D}. That is
the mapping Exercise~\ref{c2:ex:6} implements.

\begin{parts}
  \item Write out the multiplexer's truth table (four rows, one per selector combination).
  \item Derive the sum-of-products equation for \code{X} in terms of \code{A}--\code{D} and
    \code{S[1:0]} (\secref{c2:a4} works the same way for the 8-to-1 case; halve it here).
  \item Build and simulate the network in CircuitVerse, and confirm \code{X} follows each data input
    in turn as you change \code{S[1:0]}.
\end{parts}
Keep your answer: Exercise~\ref{c2:ex:6} below implements this exact circuit in VHDL, once
\secref{c2:a5} has given you the \code{process} and \code{case} statements it needs.

\exercise{A 2-to-1 mux with \code{if}}{VHDL}
\label{c2:ex:5}
\Secref{c2:a5} writes the 2:1 multiplexer using a \code{case} statement. Write the same circuit
again with an \code{if}/\code{else} instead, and confirm for yourself that the two describe identical
hardware.

Write an entity named \code{mux2} with:

\begin{booktable}{@{}p{4em}p{4.8em}p{5.2em}L@{}}
  \thead{Port} & \thead{Direction} & \thead{Type} & \thead{Description} \\
  \midrule
  \code{d0} & in & \code{std_logic} & Data input. \\
  \code{d1} & in & \code{std_logic} & Data input. \\
  \code{sel} & in & \code{std_logic} & Selector. \\
  \code{x} & out & \code{std_logic} & The selected data input. \\
\end{booktable}

\noindent The selector routes the data input whose number it names:
\begin{itemize}
  \item \code{sel = '0'} selects \code{d0}.
  \item \code{sel = '1'} selects \code{d1}.
  \item Note this is the opposite way round from \secref{c2:a5}'s \code{case} example, which selects
    its \emph{second} input \code{b} on \code{'0'}. Which input a selector value names is a
    convention, not a rule, so read it off the specification every time rather than assuming.
\end{itemize}
Implement the multiplexer using a process-based architecture:
\begin{itemize}
  \item Include \code{d0}, \code{d1}, and \code{sel} in the sensitivity list.
  \item Use an \code{if}/\code{else} statement on \code{sel}.
  \item Assign \code{x} in every branch.
  \item Do not use:
    \begin{itemize}
      \item A conditional concurrent assignment (\code{x <= d1 when sel = '1' else d0;}).
      \item A \code{with...select} statement:
        \begin{itemize}
          \item a concurrent form of \code{case} this book doesn't cover.
          \item mentioned only so you don't reach for it if you've met it elsewhere.
        \end{itemize}
    \end{itemize}
\end{itemize}

\modulefigure{mux2}

\selfcheck name your entity \code{mux2}, with ports \code{d0}, \code{d1}, \code{sel} (in) and
\code{x} (out), declared in that order; its testbench is in \file{exercises/mux2} and checks all
eight combinations of the three inputs.

\exercise{The 4-to-1 mux in VHDL}{VHDL}
\label{c2:ex:6}
Extend Exercise~\ref{c2:ex:5} into a 4:1 multiplexer named \code{mux4}.

This is the circuit you derived, drew, and simulated by hand in Exercise~\ref{c2:ex:4} above.
Re-derive the \code{case} branches yourself for four inputs and two selector bits rather than
copying the 8-to-1 example from \secref{c2:a5}.

The entity must have:

\begin{booktable}{@{}p{4em}p{4.8em}p{14.6em}L@{}}
  \thead{Port} & \thead{Direction} & \thead{Type} & \thead{Description} \\
  \midrule
  \code{d0} & in & \code{std_logic} & Data input. \\
  \code{d1} & in & \code{std_logic} & Data input. \\
  \code{d2} & in & \code{std_logic} & Data input. \\
  \code{d3} & in & \code{std_logic} & Data input. \\
  \code{sel} & in & \code{std_logic_vector(1 downto 0)} & Selector. \\
  \code{x} & out & \code{std_logic} & The selected data input. \\
\end{booktable}

\noindent As in Exercise~\ref{c2:ex:5}, the selector names its data input directly:
\begin{itemize}
  \item \code{"00"} selects \code{d0}.
  \item \code{"01"} selects \code{d1}.
  \item \code{"10"} selects \code{d2}.
  \item \code{"11"} selects \code{d3}.
  \item Again, that is the opposite mapping from \secref{c2:a5}'s 8-to-1 worked example, where
    \code{"000"} routes \code{inputs(7)}. Derive your \code{case} branches from the four lines above,
    not from that listing.
\end{itemize}
Implement the multiplexer using a process:
\begin{itemize}
  \item Include all data inputs and \code{sel} in the sensitivity list.
  \item Use a \code{case} statement on \code{sel}, one branch per selector value above.
  \item Include a \code{when others} branch.
    \begin{itemize}
      \item Choose a safe output value for select inputs containing values such as \code{'X'},
        \code{'U'}, or \code{'Z'}.
    \end{itemize}
  \item Do not replace the \code{case} statement with an \code{if}/\code{elsif} chain.
\end{itemize}

\modulefigure{mux4}

\selfcheck name your entity \code{mux4}, with ports \code{d0}, \code{d1}, \code{d2}, \code{d3},
\code{sel} (\code{std_logic_vector(1 downto 0)}) and \code{x} (out), declared in that order; its
testbench is in \file{exercises/mux4} and sweeps all four selector values against all sixteen
data-input combinations.

% ------------------------------------------------------------------------------------------
\exerciseset{Submodules and 7-Segment Displays}
\label{c2:set:4}
Exercises~\ref{c2:ex:7} and~\ref{c2:ex:8} build one design in two halves: a module that drives a
single 7-segment display, and a top level that uses two copies of it to show a hexadecimal digit on
each. Do them in order: Exercise~\ref{c2:ex:8} instantiates what you write in
Exercise~\ref{c2:ex:7}.

A 7-segment display is seven LEDs, named \code{a} through \code{g}, arranged around a figure-eight.
Light the right subset and you get a digit. This is the layout, and it is the one you need in order
to work out the codes below:

\begin{console}
      aaaa
     f    b
     f    b
      gggg
     e    c
     e    c
      dddd
\end{console}

\noindent So \code{1} lights \code{b} and \code{c} (the two on the right), and \code{7} lights
\code{a}, \code{b} and \code{c}. Comparing the codes given for those two digits is a quick way to
confirm you have the layout the right way round before you start deriving \code{A} to \code{F}.

On the DE0-CV the displays are wired \textbf{active-low}: driving a segment to \code{'0'} lights it,
and \code{'1'} turns it off. That is why the code for \code{8}, which lights every segment, is
\code{"0000000"}, and why a blank display is \code{"1111111"}.

Throughout, \code{hex(6 downto 0)} maps to segments \code{g}, \code{f}, \code{e}, \code{d},
\code{c}, \code{b}, \code{a}: bit 6 is \code{g} and bit 0 is \code{a}.

\exercise{One 7-segment digit}{VHDL, Paper}
\label{c2:ex:7}
Write a module named \code{display} that drives one 7-segment display with a hexadecimal digit.

The entity must have:

\begin{booktable}{@{}p{4em}p{4.8em}p{14.6em}L@{}}
  \thead{Port} & \thead{Direction} & \thead{Type} & \thead{Description} \\
  \midrule
  \code{number} & in & \code{std_logic_vector(3 downto 0)} & The value to show, \code{0000}
    through \code{1111}. \\
  \code{hex} & out & \code{std_logic_vector(6 downto 0)} & The segment code for that value. \\
\end{booktable}

\noindent So feeding \code{number} the value \code{0111} must put the code for \code{7} on
\code{hex}, and feeding it \code{1110} must put the code for \code{E} on \code{hex}.

Implement it using a process:
\begin{itemize}
  \item Put \code{number} in the sensitivity list.
  \item Use a \code{case} statement on \code{number}, one branch per value.
  \item Include a \code{when others} branch, and blank the display in it.
\end{itemize}
You are given the codes for \code{0} to \code{9}. Work out \code{A} to \code{F} yourself, from the
segment layout above:

\begin{vhdlcode}
constant DISPLAY_0  : std_logic_vector(6 downto 0) := "1000000";
constant DISPLAY_1  : std_logic_vector(6 downto 0) := "1111001";
constant DISPLAY_2  : std_logic_vector(6 downto 0) := "0100100";
constant DISPLAY_3  : std_logic_vector(6 downto 0) := "0110000";
constant DISPLAY_4  : std_logic_vector(6 downto 0) := "0011001";
constant DISPLAY_5  : std_logic_vector(6 downto 0) := "0010010";
constant DISPLAY_6  : std_logic_vector(6 downto 0) := "0000010";
constant DISPLAY_7  : std_logic_vector(6 downto 0) := "1111000";
constant DISPLAY_8  : std_logic_vector(6 downto 0) := "0000000";
constant DISPLAY_9  : std_logic_vector(6 downto 0) := "0010000";
-- Todo: add DISPLAY_A through DISPLAY_F here.
constant DISPLAY_OFF: std_logic_vector(6 downto 0) := "1111111";
\end{vhdlcode}

\noindent Two of the six are lower-case on a 7-segment display, because their upper-case forms would
be indistinguishable from a digit: \code{b} would read as \code{8} and \code{d} would read as
\code{0}. \code{A}, \code{C}, \code{E} and \code{F} are upper-case.

All sixteen codes are also sitting in \file{display_tb.vhd}, since a self-checking testbench cannot
check an output without knowing what it should be. So this is not a puzzle you can be locked out of.
Derive the six anyway: it takes a couple of minutes with the segment diagram, and reading them off
the testbench teaches you nothing the \code{case} statement does not. The work of this exercise is
the sixteen-branch \code{case}, not the lookup table.

\modulefigure{display}

\selfcheck name your entity \code{display}, with ports \code{number}
(\code{std_logic_vector(3 downto 0)}) and \code{hex} (out, \code{std_logic_vector(6 downto 0)}),
declared in that order; its testbench is in \file{exercises/display} and checks all sixteen values
against the full code table, naming the digit and both codes when one is wrong.

\exercise{Two digits from submodules}{VHDL}
\label{c2:ex:8}
Write a module named \code{hex_display} that shows a two-digit hexadecimal number on two displays,
using two instances of your \code{display} from Exercise~\ref{c2:ex:7}.

The entity must have:

\begin{booktable}{@{}p{4em}p{4.8em}p{14.6em}L@{}}
  \thead{Port} & \thead{Direction} & \thead{Type} & \thead{Description} \\
  \midrule
  \code{input} & in & \code{std_logic_vector(7 downto 0)} & Eight slide switches: the top four are
    the left digit, the bottom four the right. \\
  \code{hex1} & out & \code{std_logic_vector(6 downto 0)} & Segment code for the left digit. \\
  \code{hex0} & out & \code{std_logic_vector(6 downto 0)} & Segment code for the right digit. \\
\end{booktable}

\noindent When it works:
\begin{itemize}
  \item the number on \code{input(7 downto 4)} appears on \code{hex1}.
  \item the number on \code{input(3 downto 0)} appears on \code{hex0}.
\end{itemize}
Build it with submodules, following \secref{c2:a6}:
\begin{itemize}
  \item Create two instances of \code{display}, labelled \code{display1} and \code{display0}.
  \item Connect \code{input(7 downto 4)} and \code{hex1} to \code{display1}.
  \item Connect \code{input(3 downto 0)} and \code{hex0} to \code{display0}.
  \item Use positional \code{port map}, as everywhere else in this book.
  \item Write no logic of your own in \code{hex_display}. If you find yourself writing a \code{case}
    statement here, you are duplicating Exercise~\ref{c2:ex:7} rather than reusing it.
\end{itemize}
Copy your \file{display.vhd} from Exercise~\ref{c2:ex:7} into the exercise directory, since
\code{hex_display} cannot be analyzed without it. This is the first design in the book built from
more than one file, so name \file{display.vhd} first, ahead of the file that instantiates it:

\begin{shell}
ghdl -a --std=93 display.vhd hex_display.vhd hex_display_tb.vhd
ghdl -e --std=93 hex_display_tb
ghdl -r --std=93 hex_display_tb --assert-level=error --stop-time=10ms
\end{shell}

\noindent\Secref{c2:c5} explains why the order matters.

\modulefigure{hex_display}

\selfcheck name your entity \code{hex_display}, with ports \code{input}
(\code{std_logic_vector(7 downto 0)}), \code{hex1} (out) and \code{hex0} (out, both
\code{std_logic_vector(6 downto 0)}), declared in that order; its testbench is in
\file{exercises/hex_display} and sweeps all 256 input values, checking each half independently so a
design that swaps the two digits is caught rather than passing on the values where both halves
happen to match.

\exercise{A digit as gates}{Paper, CircuitVerse}
\label{c2:ex:9}
\textbf{If you have time.} Realize one digit as a gate network, the way Chapter~\ref{c1:ch} and
\secref{c2:a1} had you do it:
\begin{itemize}
  \item Take the four inputs and one segment, say \code{a}, and derive the Boolean equation for it
    from the 16-row truth table. A Karnaugh map is worth using here: most segments minimize
    substantially.
  \item Repeat for as many of the remaining six segments as you have patience for, and build the
    result in CircuitVerse.
  \item Then copy the whole gate network to drive a second display.
\end{itemize}
That copy is the point of the exercise. Duplicating a block of gates is exactly what instantiating
\code{display} a second time did in Exercise~\ref{c2:ex:8}, and it is worth having done once by hand
to see what the one line of VHDL stands for.

% ------------------------------------------------------------------------------------------
\exerciseset{Naming an Intermediate Result}
\label{c2:set:5}

\exercise{Naming a sub-expression}{VHDL, Paper}
\label{c2:ex:10}
Write an entity named \code{combo_logic} that implements the Boolean function:
\[ x = ab + c'd \]
\Secref{c2:a2} made the case that naming a sub-expression with an internal signal costs nothing and
buys readability. This exercise is where you confirm that rather than take it on trust: you write
the function twice, once as a single expression and once with the \code{c'd} term pulled out into
an internal signal named \code{y}.

A signal declared inside an architecture is a wire. It is not a variable, it does not store
anything, and naming a sub-expression with one does not add a gate, a delay, or a flip-flop to the
circuit. It only gives a piece of the network a name, so that you, reading it later, can see what
that piece is for. Both architectures below synthesize to the same gates, and the provided
testbench passes against either.

The entity must have:

\begin{narrowtable}{@{}lll@{}}
  \thead{Port} & \thead{Direction} & \thead{Type} \\
  \midrule
  \code{a} & in & \code{std_logic} \\
  \code{b} & in & \code{std_logic} \\
  \code{c} & in & \code{std_logic} \\
  \code{d} & in & \code{std_logic} \\
  \code{x} & out & \code{std_logic} \\
\end{narrowtable}

\begin{parts}
  \item Implement the entire function using a single concurrent signal assignment:
    \begin{itemize}
      \item Do not declare any internal signals.
      \item Translate each Boolean operation into its corresponding VHDL operator.
    \end{itemize}
  \item Rewrite the architecture using an internal signal named \code{y}:
    \begin{itemize}
      \item Use \code{y} to represent the term:
        \[ c'd \]
      \item Assign the final expression to \code{x}.
      \item This is the same shape as \code{gate_network} in \secref{c2:a2}, which names its
        \code{AND} term \code{ab}.
    \end{itemize}
\end{parts}
Confirm that the two architectures describe the same circuit.

\note[Tip] Two descriptions represent the same combinational circuit when they produce the same
output for every possible input combination. With four inputs:
\begin{itemize}
  \item There are 16 possible input combinations.
  \item Write the complete truth table.
  \item Verify that both architectures produce the same value of \code{x} in every row.
\end{itemize}

\modulefigure{combo_logic}

\selfcheck name your entity \code{combo_logic}, with ports \code{a}, \code{b}, \code{c}, \code{d}
(in) and \code{x} (out), declared in that order; its testbench is in \file{exercises/combo_logic}
and checks all 16 rows. It passes against either architecture, which is the point of the comparison
above.

